% concatenated from document.tex
% !TeX encoding = UTF-8
% !TeX program = xelatex
% !BIB program = biber

\documentclass[
	paper=A4,
	pagesize=auto,
	fontsize=12pt
]{scrartcl}
\usepackage{libertine}
\recalctypearea

\usepackage[british]{babel}
\usepackage[style=alphabetic,url=false,isbn=false,doi=true,eprint=false,backend=biber]{biblatex}
\AtEveryBibitem{\clearlist{language}}
\addbibresource{bibliography.bib}

\usepackage{graphicx}
\usepackage{amsmath,amssymb,amsthm}
\usepackage{csquotes}
\usepackage{bm}
\usepackage{leftidx}

\usepackage{enumitem}
\setlist[enumerate]{label = (\roman*)}

\usepackage{zref-clever}
% Load hyperref and xurl as follows in order to let arXiv break links correctly
\usepackage{hyperref}
\usepackage{xurl}
\hypersetup{colorlinks,breaklinks=true}

\AddToHook{env/proposition/begin}{%
\zcsetup{countertype={theorem=proposition}}}
\AddToHook{env/lemma/begin}{%
\zcsetup{countertype={theorem=lemma}}}
\AddToHook{env/corollary/begin}{%
\zcsetup{countertype={theorem=corollary}}}

\theoremstyle{plain}
\newtheorem{theorem}{Theorem}[section]
\newtheorem{proposition}[theorem]{Proposition}
\newtheorem{lemma}[theorem]{Lemma}
\newtheorem{corollary}[theorem]{Corollary}

\AddToHook{env/example/begin}{%
\zcsetup{countertype={theorem=example}}}
\AddToHook{env/remark/begin}{%
\zcsetup{countertype={theorem=remark}}}

\theoremstyle{remark}
\newtheorem{example}[theorem]{Example}
\newtheorem{remark}[theorem]{Remark}

\AddToHook{env/definition/begin}{%
\zcsetup{countertype={theorem=definition}}}

\theoremstyle{definition}
\newtheorem{definition}[theorem]{Definition}

\newlength{\JZHeightOfX}
\newcommand{\JZOrcidlink}[1]{
\setlength{\JZHeightOfX}{\fontcharht\font`X}
\includegraphics[height=\JZHeightOfX]{Orcidlogo.eps}
\href{https://orcid.org/#1}{#1}
}

\begin{document}
\title{Countability Properties of Weakly Compact Sets in Asymmetric Locally Convex Spaces}
\author{
Jobst Ziebell\hspace{2em}\JZOrcidlink{0000-0002-9715-6356}\\
\small{Faculty of Mathematics and Computer Science, Friedrich-Schiller-University, Jena, Germany}
}
\date{\today}
\maketitle

\begin{abstract}
A dual pair formulation for asymmetric locally convex spaces is developed that strictly generalises the ordinary vector space setting.
The concept of a polar topology carries over to the asymmetric case and some familiar results are reproduced or generalised such as the bipolar theorem and a Mackey-Arens type theorem.
The implications of weak compactness and countability properties are studied and appear intimately connected to separation properties.
An asymmetric analogue $C_a(K)$ of the well-known space $C_p(K)$ is introduced and the properties of (relatively) (countably respectively sequentially) compact subspaces are investigated.
In particular, it is shown that $C_a([0,1])$ is not angelic.
However, for Hausdorff subspaces satisfying a simple closure condition, the different compactness conditions are equivalent and imply the Fréchet-Urysohn property.
Moreover, an analogue of the Eberlein-Šmulian theorem holds in asymmetrically normed spaces.
\end{abstract}

\section{Introduction}
%
The goal of this paper is to derive an asymmetric analogue of the Eberlein-Šmulian theorem and of particular interest are the sequential properties of weakly compact sets in asymmetrically normed spaces.
A motivation to study these comes from a simple consideration in measure theory:
Suppose one has a finite measure $\mu$ on some measurable space $(X, \mathcal{A})$ and a sequence $(f_n)_{n \in \mathbb{N}}$ of real-valued $\mu$-measurable functions on $X$.
One might consider the convergence properties of the sequence
%
\begin{equation}
n \mapsto \int_X \exp \left[ f_n \right] \mathrm{d} \mu \, .
\end{equation}
%
If it converges to a finite number, one immediately knows that for large enough $N \in \mathbb{N}$, the functions $(\max\{ f_n, 0 \})_{n = N}^\infty$ are uniformly $\mu$-integrable.
Using Komlos' theorem and the convexity of the integrand, it follows that $\{ \max\{ \sum_{m = N}^n f_m / n, 0 \} : n \in \mathbb{N}_{\ge N} \}$ is relatively compact in $L^1(\mu)$.
At the same time, one knows very little about the negative parts of the $f_n$s and thus not even whether $f_n \in L^1(\mu)$ for large $n \in \mathbb{N}$.
This phenomenon is easily captured by the asymmetric seminorm $f \mapsto \int_X \max \{f, 0 \} \mathrm{d} \mu$ leading to an asymmetrically normed space.
In fact, the set $\{ f_n : n \in \mathbb{N}_{\ge N} \}$ has powerful compactness properties in the weak topology of this asymmetrically normed space.
If this set in question has sequential properties, these enable conclusions using Fatou's lemma, the dominated convergence theorem and many more.

The natural setting of the Eberlein-Šmulian theorem is the stage of dual pairs.
However, this is a concept that has - to the knowledge of the author - not yet been treated in the context of asymmetric locally convex spaces.
A considerably more general framework has been developed for locally convex cones in \cite[Chapter II.3]{src:KeimelRoth:OrderedConesAndApproximation}, but these are built on rather abstract order structures and lack a direct connection to the comparatively simple concept of a family of seminorms.
Thus, the purpose of this paper is threefold:
%
\begin{itemize}
\item Develop a simple formulation of asymmetric dual pairs similar to the one in the realm of locally convex vector spaces.
\item Investigate the interplay of countability properties and weak compactness.
\item Apply the results to the weak topology of an asymmetrically normed space.
\end{itemize}
%
In the symmetric case, it is very instructive to embed a normed space with its weak topology in a function space $C_p(K)$, that is the set of continuous functions on a compact Hausdorff space equipped with the topology of pointwise convergence.
In the asymmetric considerations, the analogous space $C_a(K)$ becomes an asymmetric version of $C_p(K)$.
Unfortunately, the angelicity of $C_p(K)$ does not survive the passage to $C_a(K)$, but very useful compactness and sequentiality properties can be derived nonetheless.

In \zcref{sec:Preliminaries}, the basic notions are introduced and some technicalities which may be unfamiliar are given.
\zcref[S]{sec:AsymmetricDualPairs} develops the notion of an asymmetric dual pair as a pairing of a vector space and a cone.
Using this formulation, polar topologies, in particular the weak and weak-$\ast$ topologies are defined and the bipolar theorem as well as a Mackey-Arens type theorem is shown.
It is also proven that every $T_0$ asymmetric locally convex topology is naturally given as a polar topology.
In \zcref{sec:WeakCompactnessAndCountability}, the implications of weak compactness and countability properties are elucidated.
These are applied to the space $C_a(K)$ in \zcref{sec:CaK} where it is shown that several topological properties of $C_p(K)$ carry over to $C_a(K)$.
Moreover, it is shown that $C_a([0,1])$ is not angelic and one can trace the failure of angelicity to the possible non-Hausdorffness of compact subsets.
Finally, the straightforward consequences for the weak topology of an asymmetrically normed space are given in \zcref{sec:WeakTopologyOnAsymmetricallyNormedSpace}.
%
\section{Preliminaries}
\label{sec:Preliminaries}
%
\subsection{Some Topological Notions}
%
A topological space is \textbf{not} implicitly assumed to be Hausdorff.
It is \textbf{regular}, if every point has a neighbourhood basis consisting of closed sets \cite[Example 15.33.c]{src:Schlechter:HandbookOfAnalysis}.
A subset $A$ of a topological space $X$ is \textbf{compact}, if every open cover of $A$ has a finite subcover and following \cite{src:KönigKuhn:AngelicSpaces}, we shall call $A$ \textbf{relatively compact} in $X$ if every open cover of $X$ has a finite subcover of $A$.
Similarly, $A$ is \textbf{countably compact}, if every countable open cover has a finite subcover and $A$ is \textbf{relatively countably compact} in $X$, if every countable open cover of $X$ has a finite subcover of $A$.
Finally, $A$ is \textbf{sequentially compact}, if every sequence in $A$ has a convergent subsequence in $A$ and is \textbf{relatively sequentially compact} in $X$ if every sequence in $A$ has a convergent subsequence in $X$.
For completeness, we state the following:
%
\begin{corollary}
$A$ is relatively compact in $X$ if and only if every net $(x_\alpha)_{\alpha \in I}$ in $A$ has a cluster point in $X$.
$A$ is relatively countably compact in $X$ if and only if every sequence $(x_n)_{n \in \mathbb{N}}$ in $A$ has a cluster point in $X$.
\end{corollary}
%\begin{proof}
%The proof follows the ideas from \cite{src:MSE:CompactnessInAngelicSpaces,src:MSE:CompactIffEveryNetHasConvergentSubnet}.
%
%$\Rightarrow$:
%Let $(x_\alpha)_{\alpha \in I}$ be a net in $A$ and suppose it has no cluster point in $X$.
%Then, for each $x \in X$, there is an open neighbourhood $U_x$ of $X$ and some $\beta_x \in I$ such that $x_\gamma \notin U_x$ for every $\gamma \in I_{\ge \beta_x}$.
%By assumption, there is a finite subset $F \subseteq X$ with $\cup_{x \in F} U_x \supseteq A$.
%Letting $\beta \in I$ denote an upper bound of $\{ \beta_x : x \in F \}$, it follows that $x_\gamma \in A \subseteq \cup_{x \in F} U_x$ for all $\gamma \in I_{\ge \beta}$, providing a contradiction.
%
%$\Leftarrow$:
%Let $(U_\alpha)_{\alpha \in I}$ be an open cover of $X$ and suppose that there is no finite subcover of $A$.
%Note that the set 
%%
%\begin{equation}
%\mathcal{U} = \left\{ \bigcup_{m = 1}^n U_{\alpha_n} \middle| n \in \mathbb{N}, \alpha_1, \dots, \alpha_n \in I \right\}
%\end{equation}
%%
%is directed under inclusion and by assumption, no $U \in \mathcal{U}$ contains all of $A$.
%Hence, for each $U \in \mathcal{U}$, we may pick some $x_U \in A \setminus U$ and consider the net $(x_U)_{U \in \mathcal{U}}$.
%Letting $x \in X$ be arbitrary, there is some $V \in \mathcal{U}$ containing $x$.
%Hence, for every $W \in \mathcal{U}$ with $W \supseteq V$, $x_W \notin V$.
%Consequently, $x$ is not a cluster point of $(x_U)_{U \in \mathcal{U}}$.
%Since $x$ was arbitrary, $(x_U)_{U \in \mathcal{U}}$ has no cluster point, contradicting the assumption.
%\end{proof}
%
Note that if $X$ is regular, then $A$ is relatively compact if and only if its closure $\bar{A}$ is compact \cite[Proposition 17.15]{src:Schlechter:HandbookOfAnalysis}.
%
\begin{definition}[{\cite{src:KönigKuhn:AngelicSpaces}}]
\label{def:angelicity}
A topological space $X$ is \textbf{angelic}, if given a relatively countably compact subset $A \subseteq X$
%
\begin{enumerate}
\item\label{def:angelicity:RelativeCompactnessEquivalence} $A$ is relatively compact,
\item for each $a \in \mathrm{cl}\, A$, there is a sequence $(a_n)_{n \in \mathbb{N}}$ in $A$ converging to $a$.
\end{enumerate} 
\end{definition}
%
A \textbf{network} in a topological space $X$ is a collection $\mathcal{N}$ of subsets of $X$ such that for every $x \in X$ and every neighbourhood $U$ of $x$, there is some $N \in \mathcal{N}$ with $x \in N \subseteq U$.
A topological space $X$ is \textbf{Lindelöf}, if every open cover has a countable subcover of $X$.
It is \textbf{hereditary Lindelöf} if the same is true for every subspace of $X$.
The following \zcref[noref]{lem:CountableNetworkImpliesHereditaryLindelöf} is well-known, but stated here for completeness.
%
\begin{lemma}
\label{lem:CountableNetworkImpliesHereditaryLindelöf}
Let $X$ be a topological space with a countable network.
Then $X$ is hereditary Lindelöf.
\end{lemma}
\begin{proof}
Denote the countable network by $\mathcal{N}$ and let $A \subseteq X$ and $(U_\alpha)_{\alpha \in I}$ be an open cover of $A$.
For every $x \in A$, there is by definition some $N_x \in \mathcal{N}$ with $x \in N_x$ and some $\alpha_x \in I$ with $N_x \subseteq U_{\alpha_x}$.
Consequently, $\cup_{x \in A} N_x$ is a cover of $A$ by the countable collection
%
\begin{equation}
\mathcal{M} = \{ N_x : x \in A \} \subseteq \mathcal{N} \, .
\end{equation}
%
Hence, there is a countable subset $B \subseteq A$ with
%
\begin{equation}
A \subseteq \bigcup_{x \in B} N_x \subseteq \bigcup_{x \in B} U_{\alpha_x} \, .
\end{equation}
\end{proof}
%
This is a variation of a well-known statement with a family of continuous functions that separates points.
%
\begin{lemma}
\label{lem:USCSeparation}
Let $X$ be a countably compact space and $F$ a countable family of real-valued upper semicontinuous functions on $X$ with the property that
%
\begin{equation}
\forall x, y \in X \exists f \in F : f \left( x \right) < f \left( y \right) \, .
\end{equation}
%
Then $X$ is compact, $T_1$ and has a countable network.
\end{lemma}
\begin{proof}
That $X$ is $T_1$ is immediately clear.
For each $f \in F$ and every $q \in \mathbb{Q}$, consider the closed set
%
\begin{equation}
C \left( f, q \right) = \left\{ x \in X : f \left( x \right) \ge q \right\}
\end{equation}
%
and let $\mathcal{N}$ be the collection of finite intersection of all such sets.
Then $\mathcal{N}$ is countable and, in fact, a network of $X$.
To see that, let $x \in X$ and $U$ be an open neighbourhood of $x$.
Clearly, $K = X \setminus U$ is countably compact and for every $y \in K$, we find some $f_y \in F$ with $f_y(y) < f_y(x)$.
Thus, there is some $q_y \in \mathbb{Q}$ such that $x \in C ( f_y, q_y )$ and $y \notin C ( f_y, q_y )$ and we may set $V_y = X \setminus C ( f_y, q_y )$.
Since $\mathcal{N}$ is countable, $\{V_y : y \in K \}$ is a countable open cover of $K$.
Hence, there are $y_1, \dots, y_n \in K$ with $\cup_{m = 1}^n V_{y_m} \supseteq K$.
Consequently,
%
\begin{equation}
x
\in
\bigcap_{m = 1}^n C ( f_{y_m}, q_{y_m} )
=
\bigcap_{m = 1}^n \left[ X \setminus V_{y_m} \right]
=
X \setminus \bigcup_{m = 1}^n V_{y_m}
\subseteq
X \setminus K
=
U \, .
\end{equation}
%
Thus, $X$ has a countable network and by \zcref{lem:CountableNetworkImpliesHereditaryLindelöf}, $X$ is Lindelöf and thus compact.
\end{proof}
%
\begin{definition}
A \textbf{quasi-uniformity} on a set $X$ is a nonempty family $\mathcal{U}$ of subsets of $X \times X$, with
%
\begin{equation}
\begin{aligned}
\forall U \in \mathcal{U} \forall x \in X &: (x, x) \in U \\
\forall U \in \mathcal{U}, V \subseteq X \times X &: U \subseteq V \implies V \in \mathcal{U} \\
\forall U, V \in \mathcal{U} &: U \cap V \in \mathcal{U} \\
\forall U \in \mathcal{U} \exists V \in \mathcal{U} &: V \circ V \subseteq U \, ,
\end{aligned}
\end{equation}
%
where $U \circ V = \{ (x, z) \in X \times X \mid \exists y \in X : (x, y) \in U \text{ and } (y, z) \in V \}$.
$\mathcal{U}$ induces a topology $\tau_{\mathcal{U}}$ on $X$ for which a subset $A \subseteq X$ is open if and only if
%
\begin{equation}
\forall x \in A \exists U \in \mathcal{U} : U[x] \subseteq A \, ,
\end{equation}
%
where $U[x] = \{ y \in X : (x, y) \in U \}$ and it can be shown that $U[x]$ is a neighbourhood of $x$.
The collection $\mathcal{U}^{-1} = \{ U^{-1} : U \in \mathcal{U} \}$ is another quasi-uniformity on $X$ called the \textbf{conjugate quasi-uniformity}, where
%
\begin{equation}
U^{-1} = \left\{ (y, x) \in X \times X : \left( x, y \right) \in U \right\} \, .
\end{equation}
%
$\mathcal{U}$ is a \textbf{uniformity} if $\mathcal{U} = \mathcal{U}^{-1}$.
Moreover, we let $\mathcal{U}_s$ (\enquote{s} for symmetric) denote the smallest uniformity containing the quasi-uniformity $\mathcal{U}$.
Given $U \in \mathcal{U}$, the \textbf{$U$-thickening} $A_U$ of a subset $A \subseteq X$ is given as
%
\begin{equation}
A
=
\left\{ y \in X \middle| \exists x \in A : \left( x, y \right) \in A \right\}
=
\bigcup_{x \in A} U \left[ x \right] \, .
\end{equation}
\end{definition}
%
\begin{lemma}
\label{lem:ClosureAsUInveseThickening}
Let $\mathcal{U}$ be a quasi-uniformity on a set $X$ equipped with the induced topology and $A \subset X$.
Then,
%
\begin{equation}
\mathrm{cl}\, A
=
\bigcap_{U \in \mathcal{U}} U^{-1} \left[A\right] \, .
\end{equation}
\end{lemma}
\begin{proof}
Let $x \in \mathrm{cl}\, A$, i.e. for every $U \in \mathcal{U}$, there is some $x_U \in A$ with $x_U \in U[x]$.
Then, $x \in U^{-1}[x_U]$.
Conversely, if $x \in \cap_{U \in \mathcal{U}} U^{-1} [A]$, for every $U \in \mathcal{U}$, there is some $x_U \in A$ with $x \in U^{-1}[x_U]$ and thus, $x_U \in U[x]$.
\end{proof}
%
A subset $A$ of a quasi-uniform space $(X, \mathcal{U})$ is \textbf{precompact} if for all $U \in \mathcal{U}$, there are $a_1, \dots, a_n \in A$ with $\cup_{m = 1}^n U[a_m] \supseteq A$.

The \textbf{pseudocharacter} of a topological space $X$ at a point $x \in X$ is countable, if there is a countable set of open sets whose intersection is equal to $\{ x \}$.
This \zcref[noref]{lem:CountabilityAndPseudocharacter} is probably known in the literature, but the author has failed to find any such reference.
%
\begin{lemma}
\label{lem:CountabilityAndPseudocharacter}
Let $X$ be a countably compact, regular topological space, $(x_n)_{n \in \mathbb{N}}$ a sequence in $X$ with some $x \in X$ as cluster point.
If the pseudocharacter of $X$ at $x$ is countable, there is a subsequence converging to $x$.
\end{lemma}
\begin{proof}
Let $(U_n)_{n \in \mathbb{N}}$ be a decreasing sequence of open sets with $\cap_{n \in \mathbb{N}} U_n = \{ x \}$.
Since $X$ is regular, there is a decreasing sequence $(V_n)_{n \in \mathbb{N}}$ of closed neighbourhoods of $x$ with $V_n \subseteq U_n$.
Now, since $x$ is a cluster point of $(x_n)_{n \in \mathbb{N}}$, there is some $n(1) \in \mathbb{N}$ with $x_{n(1)} \in V_1$.
Inductively, for every $m \in \mathbb{N}$, one finds $n(m+1)$ with
%
\begin{equation}
n \left( m + 1 \right) \ge n \left( m \right)
\qquad\text{and}\qquad
a_{n (m+1)} \in V_{m + 1} \, .
\end{equation}
%
Since $X$ is countably compact, the sequence $(x_{n(m)})_{m \in \mathbb{N}}$ has a cluster point $y \in X$.
However, since $(V_m)_{m \in \mathbb{N}}$ is decreasing, the sequence is eventually contained in every $V_m$, such that
%
\begin{equation}
y \in \bigcap_{m \in \mathbb{N}} V_m \subseteq \bigcap_{m \in \mathbb{N}} U_m = \left\{ x \right\} \, .
\end{equation}
%
Consequently, $(x_{n(m)})_{m \in \mathbb{N}}$ has precisely one cluster point, namely $x$.
Suppose now, that $(x_{n(m)})_{m \in \mathbb{N}}$ does not converge to $x$.
Then, there is a neighbourhood $U$ of $x$ and a subsequence $(y_n)_{n \in \mathbb{N}}$ lying entirely in $X \setminus U$.
By the countable compactness, $(y_n)_{n \in \mathbb{N}}$ has a cluster point $y \in X \setminus U$.
However, that cluster point is necessarily also a cluster point of $(x_{n(m)})_{m \in \mathbb{N}}$ such that $y = x$ providing a contradiction.
Hence, $(x_{n(m)})_{m \in \mathbb{N}}$ converges to $x$.
\end{proof}
%
\subsection{Asymmetric Locally Convex Spaces}
%
Vector spaces and cones are always considered over the field of real numbers.
A function $p : X \to [0, \infty)$ on a vector space or a cone $X$ is an \textbf{asymmetric seminorm}, if
%
\begin{equation}
p \left( r x \right) = r p \left( x \right)
\qquad
p \left( x + y \right) \le p \left( x \right) + p \left( y \right) \, ,
\end{equation}
%
for all $r \ge 0$ and all $x, y \in X$.
Note that this deviates from the asymmetric seminorms on cones considered in \cite[Section 2.3.3]{src:Cobzaş:FunctionalAnalysisinAsymmetricNormedSpaces}, which include a form of definiteness as well.
Since seminorms on vector spaces lack a definiteness condition, the same term will be used for similar functions on cones in this paper.
Given a vector space or a cone $X$, a family of asymmetric seminorms $\mathcal{P}$ on $X$ induces a topology $\tau_{\mathcal{P}}$ on $X$ with a basis given by sets of the form
%
\begin{equation}
U_{x, F, \epsilon}
=
\left\{ y \in X \middle| \exists \Delta \in X : y = x + \Delta \text{ and } \forall p \in F : p \left( \Delta \right) < \epsilon \right\} \, ,
\end{equation}
%
for every $x \in X$, $\epsilon > 0$ and every finite subset $F$ of $\mathcal{P}$.
A vector space $X$ equipped with this topology is an \textbf{asymmetric locally convex space} and the addition is continuous \cite[Section 1.1.7]{src:Cobzaş:FunctionalAnalysisinAsymmetricNormedSpaces}.
The topology is $T_0$ if and only if, for every $x \in X \setminus \{ 0 \}$, there is some $p \in \mathcal{P}$ with $p(x) > 0$ or $p(-x) > 0$ \cite[Proposition 1.1.63]{src:Cobzaş:FunctionalAnalysisinAsymmetricNormedSpaces}.
Moreover, every $p \in \mathcal{P}$ is clearly upper semicontinuous with respect to $\tau_{\mathcal{P}}$.
There is also an associated quasi-uniformity $\mathcal{U}$ inducing $\tau_{\mathcal{P}}$, defined as the smallest such object containing all sets of the form
%
\begin{equation}
V_{F, \epsilon}
=
\left\{ \left( x, y \right) \in X \times X \middle|
\forall p \in F : p \left( y - x \right) < \epsilon
\right\}
\end{equation}
%
for $F$ and $\epsilon$ as above \cite[Section 1.1.7]{src:Cobzaş:FunctionalAnalysisinAsymmetricNormedSpaces}.
Equivalently, just as for an abelian topological group, $\mathcal{U}$ is generated by
%
\begin{equation}
\left\{ \left( x, x + U \right) \middle| x \in X, U \subseteq X \text{ is a neighbourhood of the origin} \right\} \, .
\end{equation}
%
Consequently, we may speak of the canonical quasi-uniformity of an asymmetric locally convex space without referring to its generating seminorms.
The \textbf{conjugate} family of asymmetric seminorms is defined as $\overline{\mathcal{P}} = \{ \overline{p} : p \in \mathcal{P} \}$, where $\overline{p}(x) = p(-x)$ for all $x \in X$.
Then, $(X, \overline{\mathcal{P}})$ is the conjugate asymmetric locally convex space and the associated quasi-uniformity is given by $\overline{U} = \{ U^{-1} : U \in \mathcal{U} \}$.
An asymmetric seminorm $p$ on $X$ for which
%
\begin{equation}
\forall x \in X : p(x) = 0 \text{ and } p(-x) = 0 \implies x = 0 \, ,
\end{equation}
%
is an \textbf{asymmetric norm} \cite[Section 1.1.1]{src:Cobzaş:FunctionalAnalysisinAsymmetricNormedSpaces}.
If $\mathcal{P}$ consists of a single asymmetric norm $p$, $(X, \mathcal{P})$ is written as $(X, p)$ and is called an \textbf{asymmetrically normed space}.
To an asymmetric locally convex space $X$, we associate the \textbf{dual space} $X^*$ given as the collection of all linear and upper semi-continuous, real-valued functions on $X$.
It is clearly a cancellative cone.
Note that every locally convex space is an asymmetric locally convex space and in that setting, the dual space corresponds to the usual definition.
A subset $H \subseteq X^*$ is \textbf{equicontinuous} if for each $\epsilon > 0$, there is a neighbourhood $U \subseteq X$ of the origin, such that
%
\begin{equation}
\forall \phi \in H \forall x \in U : \phi \left( x \right) < \epsilon \, .
\end{equation}
%
Again, in the ordinary locally convex setting, this coincides with the usual definition with $\vert \phi(x) \vert < \epsilon$, because every neighbourhood of zero contains a balanced neighbourhood of zero.
\section{Asymmetric Dual Pairs}
\label{sec:AsymmetricDualPairs}
%
\begin{definition}
A tuple $(X, Y, b)$ consisting of a vector space $X$, a cone $Y$ and a bilinear (with respect to nonnegative scalars in the second variable) mapping $b : X \times Y \to \mathbb{R}$ that is \textbf{nondegenerate} in the sense that for all $x \in X$ and all $y, z \in Y$,
%
\begin{equation}
b(x, \cdot) = 0 \implies x = 0
\quad\text{and}\quad
b( \cdot , y) = b( \cdot , z)  \implies y = z,
\end{equation}
%
is an \textbf{asymmetric dual pair}.
We shall use the abbreviation $\langle \cdot, \cdot \rangle$ instead of $b$ and also simply write that $X$ and $Y$ are \textbf{in duality}.
\end{definition}
%
Obviously, every ordinary dual pair is an asymmetric dual pair and an asymmetric dual pair is an ordinary dual pair precisely, when both spaces are vector spaces.
For the remainder of the section, we fix some $X$ and $Y$ in duality.
%
\begin{corollary}
$Y$ is cancellative.
\end{corollary}
\begin{proof}
Let $a, b, c \in Y$ and suppose that $a + c = b + c$.
Then, since $\langle \cdot, \cdot \rangle$ is bilinear,
%
\begin{equation}
\left< x, a \right> 
=
\left< x, a + c \right> - \left< x, c \right>
=
\left< x, b + c \right> - \left< x, c \right>
=
\left< x, b \right> \, ,
\end{equation}
%
for every $x \in X$.
Since $\langle \cdot, \cdot \rangle$ is nondegenerate, $a = b$ follows.
\end{proof}
%
Since $Y$ is cancellative, it can be embedded in a vector space.
Specifically, we set $\widetilde{Y} = \mathrm{span}\, \{ \langle \cdot, y \rangle : y \in Y \}$ and note that $X$ and $\widetilde{Y}$ are a dual pair of vector spaces in the obvious way.
Consequently, we may consider the \textbf{annihilator} of a subset $A \subset X$ defined as
%
\begin{equation}
\mathrm{ann} A
=
\left\{ y \in \widetilde{Y} \middle| \forall x \in A : \left< x, y \right> = 0 \right\} \, .
\end{equation}
%
\begin{proposition}
\label{prop:WeakTop}
\begin{equation}
p_y : X \to [0, \infty), x \mapsto \max \left\{ \left\langle x, y \right\rangle, 0 \right\}
\end{equation}
%
is an asymmetric seminorm.
Furthermore, the collection $\mathcal{P} = \{ p_y : y \in Y \}$ turns $(X, \mathcal{P})$ into a $T_0$ asymmetric locally convex space.
If $Y$ is a vector space, then $(X, \mathcal{P})$ is a Hausdorff, locally convex space.
There is also a natural partial order $\le$ on $X$, defined as
%
\begin{equation}
x \le x' \iff \forall y \in Y : \left< x, y \right> \le \left< x', y \right>
\end{equation}
%
for all $x, x' \in X$.
It coincides with the specialization order on $(X, \mathcal{P})$, i.e. $x \le x' \iff x' \in \mathrm{cl}\, \{ x \}$.
In particular, if $(X, \mathcal{P})$ is $T_1$, $\le$ is simply the diagonal of $X \times X$.
\end{proposition}
\begin{proof}
That $p_y$ is an asymmetric seminorm is obvious, such that $(X, \mathcal{P})$ is an asymmetric locally convex space.
Suppose $x \in X$ such that for all $y \in Y$, $p_y(x) = 0$ and $p_y(-x) = 0$.
Then, $\langle x, y \rangle = 0$ for all $y \in Y$ such that $x = 0$ by assumption and thus $(X, \mathcal{P})$ is $T_0$.
If $Y$ is a vector space, it follows that $(X, \mathcal{P})$ is a $T_0$ locally convex space and thus Hausdorff.
It is easy to see that $\le$ is a partial order.
Moreover, for all $x, x' \in X$,
%
\begin{equation}
x' \in \mathrm{cl}\, \{ x \}
\iff
\forall \epsilon > 0 \forall y \in Y : p_y \left( x - x' \right) < \epsilon
\iff
\forall y \in Y : \left< x, y \right> \le \left< x', y \right> \, .
\end{equation}
\end{proof}
%
The corresponding topology on $X$ will be denoted by $\sigma(X, Y)$ in accordance with the usual notation for dual pairs.
In addition, we shall denote the conjugate topology by $\overline{\sigma}(X, Y)$ and the smallest topology containing both of these by $\sigma_s(X, Y)$.
There is of course the corresponding dual topology as well.
%
\begin{proposition}[{\cite[Proposition 2.4.1]{src:Cobzaş:FunctionalAnalysisinAsymmetricNormedSpaces}}]
\label{prop:WeakStarTop}
\begin{equation}
q_x : Y \to [0, \infty), y \mapsto \max \left\{ \left\langle x, y \right\rangle, 0 \right\}
\end{equation}
%
is an asymmetric seminorm.
The topology $\sigma(Y, X)$ generated by $\mathcal{Q} = \{ q_x : x \in X \}$ is the topology of pointwise convergence, i.e. a net $(y_\alpha)_{\alpha \in I}$ in $Y$ $\sigma(Y, X)$-converges to $y \in Y$ if and only if
%
\begin{equation}
\forall x \in X : \lim_\alpha \left< x, y_\alpha \right> = \left< x, y \right> \, .
\end{equation}
\end{proposition}
%
Consequently, $\sigma(Y, X)$ coincides with the subspace topology induced by $\sigma(\widetilde{Y}, X)$ and the annihilator of a set is necessarily $\sigma(\widetilde{Y}, X)$-closed.
The dual space is well-behaved with respect to $\sigma(X, Y)$.
%
\begin{theorem}
$(X, \sigma(X, Y))^*$ is isomorphic to $Y$ in the sense that
%
\begin{enumerate}
\item for every $\phi \in (X, \sigma(X, Y))^*$, there is a unique $y \in Y$ such that $\phi = \langle \cdot, y \rangle$,
\item for every $y \in Y$, $\langle \cdot, y \rangle \in (X, \sigma(X, Y))^*$.
\end{enumerate}
\end{theorem}
\begin{proof}
$(ii)$:
Clearly, $\langle \cdot, y \rangle \le p_y$ such that the claim follows from \cite[Proposition 2.1.11]{src:Cobzaş:FunctionalAnalysisinAsymmetricNormedSpaces}.

$(i)$:
Since $\phi$ is continuous on the asymmetrically normed space $(X, \sigma(X, Y))^*$, there are $y_1, \dots, y_n$ and some $C > 0$ with
%
\begin{equation}
\phi \left( x \right)
\le
C
\max \left\{ p_{y_1}\left( x \right), \dots, p_{y_n}\left( x \right) \right\}
=
C \max \left\{ \left< x, y_1 \right>, \dots, \left< x, y_n \right>, 0 \right\}  \, ,
\end{equation}
%
for all $x \in X$ \cite[Proposition 2.1.11]{src:Cobzaş:FunctionalAnalysisinAsymmetricNormedSpaces}.
The set $K = \{ C y_1, \dots, C y_n \}$ is clearly $\sigma(Y, X)$-compact and
%
\begin{equation}
\forall x \in X:
\sup_{y \in K} \left< x, y \right> \le 1
\implies
\phi \left( x \right) \le 1 \, .
\end{equation}
%
Hence, the existence of the sought $y \in Y$ follows from \cite[Theorem 3.8]{src:KeimelRoth:OrderedConesAndApproximation}.
The uniqueness is clear from the nondegeneracy of $\langle \cdot, \cdot \rangle$. 
\end{proof}
%
Consequently, we shall use the identification $Y \simeq (X, \sigma(X, Y))^*$ from here on.
With that in mind, it is clear that $\sigma(X, Y)$ coincides with the \textbf{weak topology} of the asymmetric locally convex space $(X, \sigma(X, Y))$ as defined in \cite[Section 2.4.6]{src:Cobzaş:FunctionalAnalysisinAsymmetricNormedSpaces}.
In particular, a net $(x_\alpha)_{\alpha \in I}$ $\sigma(X, Y)$-converges to $x \in X$ if and only if $\limsup_{\alpha} \langle x_\alpha, y \rangle \le \langle x, y \rangle$ for all $y \in Y$.
%
\begin{corollary}
\label{cor:XandXDualAreInDuality}
Let $X$ be a $T_0$ asymmetric locally convex space.
Then, $X$ and $X^*$ are in duality and $\sigma(X, X^*)$ is the weakest $T_0$ asymmetric locally convex topology on $X$ preserving the dual space $X^*$.
\end{corollary}
\begin{proof}
We just need to prove that $X$ and $X^*$ are in duality, then the rest is obvious.
If $\phi, \psi \in X^*$ with $\phi(x) = \psi(x)$ for all $x \in X$, then $\phi = \psi$ by definition.
Conversely, suppose that $x \in X$ with $\phi(x) = 0$ for all $\phi \in X^*$.
If $x \neq 0$, since $X$ is $T_0$, there is some asymmetric seminorm $p$ defining the topology with $p(x) > 0$ or $p(-x) < 0$.
Consequently, the Hahn-Banach theorem \cite[Theorem 2.2.2]{src:Cobzaş:FunctionalAnalysisinAsymmetricNormedSpaces} gives a linear functional $\psi \in X^*$ with $\psi(x) \neq 0$, providing a contradiction.
\end{proof}
%
Similarly, $\sigma(Y, X)$ coincides with the \enquote{$w^\flat$}-topology defined in \cite[Section 2.4.6]{src:Cobzaş:FunctionalAnalysisinAsymmetricNormedSpaces} and thus deserves the name \textbf{weak-$\ast$ topology}.
Call a set $A \subseteq X$ \textbf{$\sigma(X,Y)$-bounded} respectively $B \subseteq Y$ \textbf{$\sigma(Y,X)$-bounded} if
%
\begin{equation}
\forall y \in Y : \sup_{x \in A} \left< x, y \right> < \infty
\quad\text{respectively}\quad
\forall x \in X : \sup_{y \in B} \left< x, y \right> < \infty \, .
\end{equation}
%
It is easy to see that this boundedness coincides with the boundedness in ordinary dual pairs.
%
\begin{proposition}
Let $A \subseteq X$ and $B \subseteq Y$.
Then the function
%
\begin{equation}
p_B : X \to [0, \infty], x \mapsto \sup_{y \in B} p_y \left( x \right)
\end{equation}
%
is an asymmetric seminorm if and only if $B$ is $\sigma(Y, X)$-bounded.
Similarly,
%
\begin{equation}
q_A : Y \to [0, \infty], y \mapsto \sup_{x \in A} q_x \left( y \right)
\end{equation}
%
is an asymmetric seminorm if and only if $A$ is $\sigma(X, Y)$-bounded.
\end{proposition}
%
Consequently, we can consider \textbf{polar topologies} on $X$ respectively $Y$ generated by suitable sets on $Y$ respectively $X$, just as in the vector space setting.
Letting $\mathcal{G}$ denote an arbitrary collection of subsets of $Y$ respectively $X$, the corresponding topology induced by the asymmetric seminorms $\{ p_G : G \in \mathcal{G} \}$ respectively $\{ q_G : G \in \mathcal{G} \}$ will be called the \textbf{$\mathcal{G}$-topology} on $X$ respectively $Y$.
In particular $\sigma(X, Y)$ and $\sigma(Y, X)$ are - as usual - generated by the collection of all single point sets in $Y$ respectively in $X$.
%
\begin{proposition}
\label{prop:PolarProperties}
Let $A \subset X$ and $B \subseteq Y$.
Define the \textbf{polar} $A^\circ$ and the \textbf{prepolar} $\leftidx{^\circ}{B}{}$ as
%
\begin{equation}
\begin{aligned}
A^\circ = \left\{ y \in Y : q_A \left( y \right) \le 1 \right\} &= \left\{ y \in Y : \sup_{x \in A} \left< x, y \right> \le 1 \right\} \, \\
\leftidx{^\circ}{B}{} = \left\{ x \in Y : p_B \left( x \right) \le 1 \right\} &= \left\{ x \in X : \sup_{y \in B} \left< x, y \right> \le 1 \right\} \, .
\end{aligned}
\end{equation}
%
Then $A^\circ$ is $\sigma(Y, X)$-closed and convex.
Moreover, it is \textbf{nonnegatively balanced} in the sense that $y \in A^\circ$ implies $t y \in A^\circ$ for all $t \in [0, 1]$.
Similarly, $\leftidx{^\circ}{B}{}$ is $\overline{\sigma}(X, Y)$-closed, convex and nonnegatively balanced.
\end{proposition}
\begin{proof}
That $A^\circ$ and $\leftidx{^\circ}{B}{}$ are convex and nonnegatively balanced is trivial.
Moreover, $A^\circ$ is $\sigma(Y, X)$-closed by \cite[Proposition 2.4.29]{src:Cobzaş:FunctionalAnalysisinAsymmetricNormedSpaces}.
To see that $\leftidx{^\circ}{B}{}$ is $\overline{\sigma}(X, Y)$-closed, let $(x_\alpha)_{\alpha \in I}$ be a net in $\leftidx{^\circ}{B}{}$ that $\overline{\sigma}(X, Y)$-converges to some $x \in X$.
Then,
%
\begin{equation}
\left< x, y \right>
\le
\liminf_{\alpha} \left< x_\alpha, y \right>
\le
1 \, ,
\end{equation}
%
for all $y \in Y$ such that $x \in \leftidx{^\circ}{B}{}$.
\end{proof}

\begin{theorem}[Bipolar Theorem]
\label{thm:bipolar}
Let $A \subset X$ and $B \subseteq Y$.
Letting $\mathcal{N}$ denote a neighbourhood basis of the origin in $\sigma(X, Y)$,
%
\begin{align}
\label{eq:BipolarWeakConjugateClosureRep}
\leftidx{^\circ}{\left( A^\circ \right)}{}
&=
\mathrm{cl}_{\overline{\sigma}(Y,X)} \mathrm{conv}\, \left[ A \cup \{ 0 \} \right]
=
\bigcap_{N \in \mathcal{N}} \left[ N + \mathrm{conv}\, \left( A \cup \{ 0 \} \right) \right]
\\
\left( \leftidx{^\circ}{B}{} \right)^\circ
&=
\mathrm{cl}_{\sigma(Y,X)} \mathrm{conv}\, \left[ B \cup \{ 0 \} \right] \, .
\end{align}
\end{theorem}
\begin{proof}
We begin, by showing the bipolar theorem for $A$.
Trivially, $0 \in \leftidx{^\circ}{( A^\circ )}{}$ and $\langle x, y \rangle \le 1$ for all $x \in A$, $y \in A^\circ$, such that $A \subseteq \leftidx{^\circ}{( A^\circ )}{}$.
By \zcref{prop:PolarProperties}, it follows that
%
\begin{equation}
\mathrm{cl}_{\overline{\sigma}(Y,X)} \mathrm{conv}\, \left( A \cup \{ 0 \} \right)
\subseteq
\leftidx{^\circ}{\left( A^\circ \right)}{} \, .
\end{equation}
%
If strict inequality holds, there is some $x \in \leftidx{^\circ}{( A^\circ )}{}$ that is not contained in $\mathrm{cl}_{\overline{\sigma}(Y,X)} \mathrm{conv}\, ( A \cup \{ 0 \} )$.
Consequently, we may apply the Hahn-Banach theorem \cite[Theorem 2.2.9]{src:Cobzaş:FunctionalAnalysisinAsymmetricNormedSpaces} and find a $\phi \in (X, \overline{\sigma}(X, Y))^*$ with
%
\begin{equation}
\phi \left( x \right)
<
\inf \phi \left( \mathrm{cl}_{\overline{\sigma}(Y,X)} \mathrm{conv}\, ( A \cup \{ 0 \} ) \right) \, .
\end{equation}
%
By definition, $\phi$ is $\overline{\sigma}(X, Y)$-upper semicontinuous, and thus $\psi = -\phi \in Y$ with
%
\begin{equation}
\sup \psi \left( \mathrm{cl}_{\overline{\sigma}(Y,X)} \mathrm{conv}\, ( A \cup \{ 0 \} ) \right)
<
\psi \left( x \right) \, .
\end{equation}
%
Since the left-hand side contains the origin, $\psi(x) > 0$ and by the strict inequality, there is some $r > 1$ with $r \psi / \psi(x) \in A^\circ$.
However, this contradicts $x \in \leftidx{^\circ}{( A^\circ )}{}$.
Consequently,
%
\begin{equation}
\leftidx{^\circ}{\left( A^\circ \right)}{}
=
\mathrm{cl}_{\overline{\sigma}(Y,X)} \mathrm{conv}\, \left( A \cup \{ 0 \} \right) \, .
\end{equation}
%
The rest of \zcref{eq:BipolarWeakConjugateClosureRep} follows from \zcref{lem:ClosureAsUInveseThickening}.

For $B$, the claim is obvious by recalling that $\sigma(Y, X)$ coincides with the subspace topology induced by $\sigma(\widetilde{Y}, X)$ and applying the bipolar theorem for dual pairs of vector spaces. 
\end{proof}
%
Let $\tau(X,Y)$ denote the \textbf{Mackey topology} on $X$, i.e. the $\mathcal{G}$-topology induced by
%
\begin{equation}
\mathcal{G}
=
\left\{ B \subseteq Y : B \text{ is convex and $\sigma(Y,X)$-compact} \right\} \, .
\end{equation}
%
Note that, in any Hausdorff, topological vector space, the convex, balanced hull of a convex, compact set is a compact disk \cite[p.66, 10.2]{src:SchaeferWolff:TVS}.
Hence, in the vector space setting, $\tau(X,Y)$ coincides with the ordinary Mackey topology.
%
\begin{lemma}
\label{lem:MackeyTopIsFiner}
Let $X$ be a vector space with a $T_0$ asymmetric locally convex topology $\alpha$.
Then, $\alpha \subseteq \tau(X, (X, \alpha)^*)$.
\end{lemma}
\begin{proof}
Set  $Y = (X, \alpha)^*$ for brevity and let $\mathcal{P}$ denote a set of asymmetric seminorms on $X$ generating its topology.
Furthermore, let $\beta$ denote the asymmetric locally convex topology induced by the conjugate family $\overline{\mathcal{P}}$ of seminorms.
For every $p \in \mathcal{P}$, the set $p^{-1}([0, 1] )$ is $\beta$-closed \cite[Proposition 1.1.59]{src:Cobzaş:FunctionalAnalysisinAsymmetricNormedSpaces} and moreover, the polar $K = p^{-1}([0, 1])^\circ$ in the pairing $(X, Y)$ is a convex, $\sigma(Y, X)$-compact set \cite[Theorem 2.4.30]{src:Cobzaş:FunctionalAnalysisinAsymmetricNormedSpaces}.
Furthermore, by \zcref{thm:bipolar},
%
\begin{equation}
\forall x \in X : p_K \left( x \right) \le 1
\iff
p \left( x \right) \le 1 \, .
\end{equation}
%
Hence, $p = p_K$, i.e. $p$ is an asymmetric seminorm defining the Mackey topology.
\end{proof}
%
\begin{theorem}
Let $X$ and $Y$ be in duality.
Then, $\tau(X, Y)^*$ is the finest $T_0$ asymmetric locally convex topology on $X$ with dual space $Y$.
\end{theorem}
\begin{proof}
We follow the proof given in \cite[p. 131, Theorem 3.2]{src:SchaeferWolff:TVS} for the vector space setting.
It is clear, that $\tau(X,Y)$ is finer than $\sigma(X,Y)$, such that
%
\begin{equation}
Y \simeq (X, \sigma(X, Y))^* \subseteq (X, \tau(X, Y))^* =: Z \, .
\end{equation}
%
For $\phi \in Z$ there exist $C > 0$ and $\sigma(Y,X)$-compact, convex sets $B_1, \dots B_n \subseteq Y$ \cite[Proposition 2.1.11]{src:Cobzaş:FunctionalAnalysisinAsymmetricNormedSpaces} with
%
\begin{equation}
\phi
\le
C \max \left\{ p_{B_1}, \dots, p_{B_n} \right\} \, .
\end{equation}
%
Recall that $\sigma(Y, X)$ coincides with the subspace topology induced by $\sigma(\widetilde{Y}, X)$ and that the convex hull of finitely many convex compact sets in a topological vector space is compact.
Then,
%
\begin{equation}
B = \mathrm{conv}\, \left[ \left\{ 0 \right\} \cup C \bigcup_{m = 1}^n B_m \right]
\end{equation}
%
is also convex and $\sigma(Y, X)$-compact, i.e. $B \in \mathcal{G}$.
Clearly, $\phi \le p_B$ and thus,
%
\begin{equation}
\leftidx{^\circ}{\left\{ \phi \right\}}{}
=
\left\{ x \in X : \phi \left( x \right) \le 1 \right\}
\supseteq
\left\{ x \in X : p_B \left( x \right) \le 1 \right\}
=
\leftidx{^\circ}{B}{} \, ,
\end{equation}
%
where the prepolars are considered in the dual pair $(X, Z)$.
Taking polars in $(X, Z)$ on both sides, we thus have, $( \leftidx{^\circ}{\{ \phi \}}{} )^\circ \subseteq ( \leftidx{^\circ}{B}{} )^\circ$.
Equipping both $Y$ and $Z$ with their respective weak-$\ast$ topologies, it is clear that the canonical injection $Y \to Z$ is linear and continuous.
Since $B$ is convex and compact in $Y$, it is thus also convex and compact in $Z$.
Hence, by the bipolar theorem,
%
\begin{equation}
\phi \subseteq \left( \leftidx{^\circ}{\left\{ \phi \right\}}{} \right)^\circ \subseteq \left( \leftidx{^\circ}{B}{} \right)^\circ = B \, .
\end{equation}
%
Consequently, $\phi \in Y$ and $Y = Z$.
That the Mackey topology is the finest asymmetric locally convex topology with dual space $Y$ follows immediately from \zcref{lem:MackeyTopIsFiner}.
\end{proof}
%
The following results constitute the asymmetric analogues of the Mackey-Arens theorem.
%
\begin{corollary}
Let $X$ be a $T_0$ asymmetric locally convex space $X$ with topology $\alpha$.
Then, $\sigma(X, X^*) \subseteq \alpha \subseteq \tau(X, X^*)$.
\end{corollary}
%
\begin{corollary}
Let $\alpha$ be a $T_0$ asymmetric locally convex topology on $X$ with $\sigma(X, Y) \subseteq \alpha \subseteq \tau(X, Y)$.
Then $(X, \alpha)^* = Y$.
\end{corollary}
%
Finally, every $T_0$ asymmetric locally convex topology is naturally a $\mathcal{G}$-topology.
%
\begin{proposition}
\label{prop:PolarOfNeighbourhoodIsEquicontinuous}
Let $X$ be an asymmetric locally convex space and $V \subseteq X$ a neighbourhood of the origin.
Then $V^\circ$ is equicontinuous.
\end{proposition}
%\begin{proof}
%Let $\delta > 0$ and $F$ denote a finite subfamily of asymmetric seminorms on $X$ inducing its topology with $U_{0, F, \delta} \subseteq V$.
%Now, for every $\epsilon > 0$, just note that $\epsilon U_{0, F, \delta} / 2$ is an open neighbourhood of zero and
%%
%\begin{equation}
%\forall \phi \in U_{0, F, \delta}^\circ \forall x \in \epsilon U_{0, F, \delta} / 2 :
%\phi \left( x \right)
%\le
%\frac{\epsilon}{2}
%<
%\epsilon \, .
%\end{equation}
%%
%Since $V^\circ \subseteq U_{0, F, \delta}^\circ$, the claim follows.
%\end{proof}
%
\begin{theorem}
Let $X$ be a $T_0$ asymmetric locally convex space and set
%
\begin{equation}
\mathcal{G} = \left\{
H \subseteq X^* : \text{$H$ is equicontinuous}
\right\} \, .
\end{equation}
%
Then the $\mathcal{G}$-topology on $X$ coincides with its original topology.
\end{theorem}
\begin{proof}
First, we need to show that every $H \in \mathcal{G}$ is $\sigma(X^*, X)$-bounded.
Let $\epsilon > 0$ and $U \subseteq X$ be a neighbourhood of the origin with $H(U) \subseteq ( -\infty, \epsilon )$.
For every $x \in X$, there is clearly some $t > 0$ with $t x \in U$ such that
%
\begin{equation}
\sup_{\phi \in H} \phi \left( x \right)
=
\frac{1}{t} \sup_{\phi \in H} \phi \left( t x \right)
<
\frac{\epsilon}{t} \, .
\end{equation}
%
Hence, $H$ is $\sigma(X^*, X)$-bounded and $\mathcal{G}$ indeed defines an asymmetric locally convex topology $\tau'$ on $X$.
Letting $\tau$ denote the original topology on $X$ and $\overline{\tau}$ the conjugate topology, $\tau$ has a neighbourhood basis $\mathcal{N}$ at the origin consisting of convex, $\overline{\tau}$-closed sets \cite[Section 1.1.7]{src:Cobzaş:FunctionalAnalysisinAsymmetricNormedSpaces}.
Moreover, for every $U \in \mathcal{N}$, the polar $U^\circ$ is equicontinuous by \zcref{prop:PolarOfNeighbourhoodIsEquicontinuous} such that $p_{U^\circ}$ is an upper semicontinuous asymmetric seminorm on $\tau'$.
By the bipolar theorem, 
%
\begin{equation}
p_{U^\circ}^{-1}([0, 1))
\subseteq
p_{U^\circ}^{-1}([0, 1])
=
U \, .
\end{equation}
%
Since $p_{U^\circ}^{-1}([0, 1))$ is a $0$-neighbourhood in $\tau'$, so is $U$ and thus $\tau \subseteq \tau'$.

Conversely, let $U \subseteq X$ be a $0$-neighbourhood in $\tau'$.
Then, there are equicontinuous subsets $H_1, \dots, H_n \subseteq Y$ and some $\epsilon > 0$ such that
%
\begin{equation}
\cap_{m = 1}^n p_{H_m}^{-1} \left( [0, \epsilon ) \right)
\subseteq
U \, .
\end{equation}
%
Clearly, the set $H = \cup_{m = 1}^n H_m$ is equicontinuous as well and $p_H^{-1}([0, \epsilon)) \subseteq U$.
By definition, there is a $\tau$-neighbourhood $V$ of the origin with $V \subseteq p_H^{-1}([0, \epsilon / 2))$ and thus $U$ is a $0$-neighbourhood in $\tau$, i.e. $\tau' \subseteq \tau$.
\end{proof}
%
\section{Countability Conditions and Weak Compactness}
\label{sec:WeakCompactnessAndCountability}
%
By taking linear (conic) combinations with rational coefficients, the following is obvious.
%
\begin{proposition}
Let $X$ and $Y$ be in duality and equip both with any $\mathcal{G}$-topologies.
If $A \subseteq X$ respectively $B \subseteq Y$ is separable, then $\mathrm{span}\, A$ respectively $\mathrm{cone}\, B$ is separable.
\end{proposition}
%
For the remainder of this section, $X$ and $Y$ are taken to be in duality and are assumed to carry their respective weak topologies.
The following \zcref[noref]{prop:VectorDualitySeparableVersusCountablyCompact} appears to be folklore.
%
\begin{proposition}
\label{prop:VectorDualitySeparableVersusCountablyCompact}
Suppose $X$ is separable and $Y$ is a vector space.
Then every countably compact subset of $Y$ is compact and metrisable.
\end{proposition}
\begin{proof}
Let $A \subseteq Y$ be countably compact and $(x_n)_{n \in \mathbb{N}}$ a dense sequence in $X$.
For each $n \in \mathbb{N}$, consider the function
%
\begin{equation}
f_n : Y \to \mathbb{R}, y \mapsto \left< x_n, y \right> \, .
\end{equation}
%
Dually, we may consider the functions $g_y : X \to \mathbb{R}, x \mapsto \langle x, y \rangle$ for all $y \in Y$.
By definition of $\sigma(X, Y)$, each $g_y$ is upper semicontinuous.
Since $Y$ is a vector space, so is $g_{-y} = - g_y$ such $g_y$ is, in fact, continuous.
Hence, it is immediate that $G = F \cup (-F)$ with $F = \{ f_n : n \in \mathbb{N} \}$ satisfies the conditions of \zcref{lem:USCSeparation}.
Since \zcref{prop:WeakStarTop} implies that $S$ is Hausdorff, the claim follows from \cite[Theorem 3.1.19]{src:Engelking:GeneralTopology}.
\end{proof}
%
In the asymmetric setting separability is not enough, but a countable basis suffices.
%
\begin{lemma}
\label{lem:CountableBasisImpliesMetrisableCompactaInDual}
If $X = \mathrm{span}\, E$ for some countable set $E \subseteq X$, every countably compact subset of $Y$ is compact and metrisable.
\end{lemma}
\begin{proof}
Let $S \subseteq Y$ be countably compact and $(x_n)_{n \in \mathbb{N}}$ an enumeration of $E$.
By the non-degeneracy, it is clear that whenever there are $y, z \in Y$ with $\langle x_n, y \rangle = \langle x_n, z \rangle$ for all $n \in \mathbb{N}$, then $y = z$.
Hence, the set $F = \{ f_n : n \in \mathbb{N} \}$ with
%
\begin{equation}
f_n : Y \to \mathbb{R}, y \mapsto \left< x_n, y \right> \, ,
\end{equation}
%
separates points in $Y$.
Since each $f_n$ is continuous by \zcref{prop:WeakStarTop}, the set $G = F \cup (-F)$ satisfies the conditions of \zcref{lem:USCSeparation}.
Hence, $S$ is compact and has a countable network
\zcref[S]{prop:WeakStarTop} and \cite[Theorem 3.1.19]{src:Engelking:GeneralTopology} together imply the metrisability.
\end{proof}
%
In the reverse setting, separability suffices, but requires additional separation assumptions.
%
\begin{theorem}
\label{thm:SeparableDualImpliesQuasimetrisableComplicatedSetsBase}
If $Y$ is separable, every subset $S \subseteq X$ that is $T_1$ and countably compact is compact and has a countable network.
\end{theorem}
\begin{proof}
First, note that for every $y \in Y$, the function
%
\begin{equation}
f_y : X \to \mathbb{R}, x \mapsto \left< x, y \right>
\end{equation}
%
is upper semicontinuous.
Now, let $x_1, x_2 \in S$ be distinct.
By assumption, each has a neighbourhood not containing the other, such that there are $y_1, \dots, y_n \in Y$ and some $\epsilon > 0$ with
%
\begin{equation}
\max_{m \in \{1, \dots, n \}} \left< x_2 - x_1, y_m \right>
>
\epsilon \, .
\end{equation}
%
In particular, there is some $y' \in Y$ with $\langle x_1, y' \rangle < \langle x_2, y' \rangle$.
Letting $D \subseteq Y$ denote a countable $\sigma(Y, X)$-dense set, we see from \zcref{prop:WeakStarTop}, that there is some $y \in D$ with $\langle x_1, y \rangle < \langle x_2, y \rangle$.
Consequently, \zcref{lem:USCSeparation} applies.
\end{proof}
%
\begin{remark}
It is unknown to the author, if anything useful can be said if $S$ is considered without the $T_1$ assumption.
However, recall that $X$ is $T_0$ by \zcref{prop:WeakTop}.
\end{remark}
%
\begin{remark}
\label{rem:HausdorffSubsetAndSeparableDualImpliesMetrisability}
If the subset in question is Hausdorff, then metrisability follows \cite[Theorem 3.1.19]{src:Engelking:GeneralTopology}.
\end{remark}
%
If $X$ is Lindelöf, every closed countably compact set is clearly compact.
Curiously, given the Lindelöf property of $\sigma_s(X, Y)$, the failure of a countably compact set to be compact is connected to the non-closedness in the symmetric topology.
%
\begin{lemma}
\label{lem:CountableCompactnessImpliesRelativeCompactnessInSymmetricClosure}
If $(X, \sigma_s(X, Y))$ is Lindelöf, every countably compact subset of $X$ is precompact.
Moreover, every countably compact $A \subseteq X$ is relatively compact in $B = \mathrm{cl}_{\sigma_s(X, Y)} A$ with the subspace topology inherited from $X$.
\end{lemma}
\begin{proof}
Let $\mathcal{U}$ denote the canonical quasi-uniformity on $(X, \sigma(X, Y))$ and let $U \in \mathcal{U}$.
Then, $U[A] \supseteq B$ by \zcref{lem:ClosureAsUInveseThickening}.
Since $B$ is $\sigma_s(X, Y)$-closed, $B$ is $\sigma_s(X, Y)$-Lindelöf and thus also $\sigma(X, Y)$-Lindelöf.
Consequently, there is a countable $S \subseteq A$ with $\cup_{a \in S} U[a] \supseteq B \supseteq A$.
Because $A$ is countably compact, there is now a finite $F \subseteq S$ with $\cup_{a \in F} U[a] \supseteq A$.
Similarly, since $B$ is Lindelöf, every open cover of $B$ has a countable subcover and thus a finite subcover of $A$.
\end{proof}
%
\begin{theorem}
\label{thm:LindelöfCompactnessCondition}
If $(X, \sigma_s(X, Y))$ is Lindelöf, a countably compact subset $A \subseteq X$ is compact if and only if
%
\begin{equation}
\label{eq:SymmetricClosureContainedInDOwnwardClosure}
\mathrm{cl}_{\sigma_s(X, Y)} A \subseteq \downarrow A \, ,
\end{equation}
%
where $\downarrow A$ denotes the downward closure of $A$ in the partial order introduced in \zcref{prop:WeakTop}.
\end{theorem}
\begin{proof}
$\Rightarrow$:
Suppose that \zcref{eq:SymmetricClosureContainedInDOwnwardClosure} holds and let $(x_\alpha)_{\alpha \in I}$ be a net in $A$.
By \zcref{lem:CountableCompactnessImpliesRelativeCompactnessInSymmetricClosure}, there is a cluster point $x \in \mathrm{cl}_{\sigma_s(X, Y)} A$ and thus $x \in \downarrow A$.
Hence, there is some $x' \in A$ with $x \in U$ for every open neighbourhood $U$ of $x'$.
Consequently, $x'$ is also a cluster point.

$\Leftarrow$:
Suppose that $A$ is compact and let $(x_\alpha)_{\alpha \in I}$ be a net in $A$ $\sigma_s(X, Y)$-converging to some $x \in \mathrm{cl}_{\sigma_s(X, Y)}$.
By compactness, there is a cluster point $x' \in A$ and it is obvious that $x \le x'$ and thus $x \in \downarrow A$.
\end{proof}
%
This raises the question, whether the closure in the symmetric topology preserves countable compactness.
%
\begin{theorem}
\label{thm:LindelöfAndSymmetricClosureImpliesCompactness}
If $(X, \sigma_s(X, Y))$ is Lindelöf, the $\sigma_s(X, Y)$-closure of every countably compact subset of $X$ is compact.
\end{theorem}
\begin{proof}
Let $A \subseteq X$ be countably compact and let $(b_n)_{n \in \mathbb{N}}$ be a sequence in $B = \mathrm{cl}_{\sigma_s(X, Y)} A$.
Set
%
\begin{equation}
\mathcal{V} = \left\{ U \in \mathcal{U}_s : U = U^{-1} \right\} \, .
\end{equation}
%
For every $V \in \mathcal{V}$ and every $m \in \mathbb{N}$, pick some $a^V_m \in V[b_n]$.
Then $(a^V_m)_{V \in \mathcal{V}, m \in \mathbb{N}}$ becomes a net under the partial order
%
\begin{equation}
(V, m) \le (W, n)
\iff
V \supseteq W \text{ and } m \le n \, ,
\end{equation}
%
for all $V, W \in \mathcal{V}$ and all $m, n \in \mathbb{N}$.
By \zcref{lem:CountableCompactnessImpliesRelativeCompactnessInSymmetricClosure}, there is a cluster point $b \in B$, i.e.
%
\begin{equation}
\forall U \in \mathcal{U} \forall V \in \mathcal{V} \forall m \in \mathbb{N} \exists W \in \mathcal{V}, n \in \mathbb{N} :
V \supseteq W, m \le n \text{ and } a^W_n \in U \left[ b \right] \, .
\end{equation}
%
Since $W = W^{-1}$, we have
%
\begin{equation}
b_n \in W \left[ a^W_n \right] \subseteq W \left[ U \left[ b \right] \right] \, .
\end{equation}
%
By setting $V = U \cap U^{-1}$, we thus have $b_n \in (U \circ U)[b]$.
It follows that $b$ is a cluster point of the sequence $(b_n)_{n \in \mathbb{N}}$.
Hence, $B$ is countably compact and thus compact.
\end{proof}
%
With suitable countability conditions one does not need to pass to the symmetric closure.
%
\begin{theorem}
Let $X$ and $Y$ be in duality and consider both with their respective weak topologies.
If $(X, \sigma_s(X, Y))$ has countable tightness, is Lindelöf and $Y$ is separable, every countably compact subset of $X$ is compact.
\end{theorem}
\begin{proof}
Let $A \subseteq X$ be countably compact, set $B = \mathrm{cl}_{\sigma_s(X, Y)} A$ with the subspace topology of $X$ and let $b \in B$ be arbitrary.
Since $(X, \sigma_s(X, Y))$ has countable tightness, there is a countable set $S \subseteq A$ with $b \in \mathrm{cl}_{\sigma_s(X, Y)} S$.
Now, pick a dense sequence $(y_n)_{n \in \mathbb{N}}$ of $Y$ and note that the map
%
\begin{equation}
\Phi : (X, \sigma_s(X, Y)) \to \mathbb{R}^{\mathbb{N}}, x \mapsto \left( \left< x, y_n \right> \right)_{n \in \mathbb{N}}
\end{equation}
%
is continuous.
Moreover, since the range is metrisable, there is a sequence $(a'_n)_{n \in \mathbb{N}}$ in $S$ with $\lim_{n \to \infty} \Phi(a'_n) = \Phi(b)$.
Consequently, $\lim_{n \to \infty} \langle a'_n, y_m \rangle = \langle b, y_m \rangle$ for all $m \in \mathbb{N}$.
Since $A$ is countably compact, there is a subnet $(a''_\beta)_{\beta \in J}$ converging to some $c \in A$.
Hence,
%
\begin{equation}
\left< b, y_m \right>
=
\lim_{\beta} \left< a''_\beta, y_m \right>
\le
\left< c, y_m \right> \, ,
\end{equation}
%
for all $m \in \mathbb{N}$.
By \zcref{prop:WeakStarTop}, $b \le c$, i.e. $B \subseteq \downarrow A$ and \zcref{thm:LindelöfCompactnessCondition} applies.
\end{proof}
%
Finally, a simple criterion for countable tightness is identical to the vector space setting.
%
\begin{proposition}
\label{prop:DualIsSigmaCompactImpliesCountableTightness}
Let $X$ and $Y$ be in duality and $Y = \cup_{n \in \mathbb{N}} M_n$ for some sequence $(M_n)_{n \in \mathbb{N}}$ of $\sigma(Y, X)$-compact subsets of $Y$.
Then $\sigma(X, Y)$ has countable tightness.
\end{proposition}
\begin{proof}
We mirror the proof given in \cite[Theorem 3.54]{src:Fabian:BanachSpaceTheory} for the symmetric setting.
Let $A \subseteq X$ be nonempty, $\bar{a} \in \mathrm{cl}_{\sigma(X,Y)} A$ and without loss of generality, suppose that $M_n \subseteq M_{n+1}$ for each $n \in \mathbb{N}$.
Now, for every $m, n, k \in \mathbb{N}$ and every $n$-tuple $T = (y_1, \dots, y_n) \in M_m^n$, pick some $a_T \in A$ with
%
\begin{equation}
a_T \in \left \{ x \in X \middle| \forall i \in \left \{ 1, \dots, n \right \} : \left< x - \bar{a}, y_i \right> < \frac{1}{k} \right\} \, .
\end{equation}
%
The set $V_T = \{ y \in M_m : \langle a_T - \bar{a}, y \rangle < 1/k \}$ is open in $M_m$ with the subspace topology induced by $\sigma(Y, X)$.
Furthermore, $y_1, \dots, y_n \in V_T$, such that $V_T^n$ is an open neighbourhood of $T \in M_m^n$.
Proceeding similarly for every $n$-tuple in $M_m$, we obtain an open cover of $M_m^n$ which has a finite subcover indexed by a finite family $T_1, \dots, T_j \in M_m^n$.
Setting $A_{m,n,k} = \{ a_{T_1}, \dots, a_{T_j} \}$ and $B = \cup_{m,n,k \in \mathbb{N}} A_{m,n,k}$ we have $\bar{a} \in \mathrm{cl}_{\sigma(X,Y)} B$ by construction.
\end{proof}
%
%
\section{The space \texorpdfstring{$C_a(K)$}{Cₐ(K)}}
\label{sec:CaK}
%
There is extensive literature on the space of continuous functions on a given topological space with the topology of pointwise convergence (see e.g. \cite{src:Arkhangelskii1:TopologicalFunctionSpaces}).
In this section, some basic properties of an asymmetric version thereof are analysed.
Because this appears to be the first study of such spaces, no attempt is being made to extend the very general results known in the literature.
Instead, the aim is to uncover some easily accessible properties that do or do not carry over to the asymmetric setting.
Hopefully, this can provide some intuition and a basis for more general studies.

Let $K$ denote a compact Hausdorff space and $C(K)$ the linear space of real-valued continuous functions on $K$.
Moreover, for each $x \in K$, let
%
\begin{equation}
\delta_x : C(K) \to \mathbb{R}, f \mapsto f \left( x \right) \, .
\end{equation}
%
Set $\widehat{K} = \{ \delta_x : x \in K \}$, $L = \mathrm{cone}\, \widehat{K}$ and note that $C(K)$ and $L$ are in duality in the obvious way.
Let $C_a(K)$ denote the space $C(K)$ equipped with the $\sigma(C(K), L)$-topology (\enquote{a} for \enquote{asymmetric}).
This topology is similar to the one of $C_p(K)$, i.e. the topology of pointwise convergence.
%
\begin{corollary}
\label{cor:CaKConvergence}
A net $(f_\alpha)_{\alpha \in I}$ in $C_a(K)$ converges to $f \in C_a(K)$ if and only if for all $x \in K$,
%
\begin{equation}
\limsup_\alpha f_\alpha \left( x \right) \le f \left( x \right) \, .
\end{equation}
\end{corollary}
%
\begin{corollary}
$C_a(K)$ is hyperconnected and ultraconnected (see e.g. \cite[p. 29]{src:SteenSeebach:CounterexamplesInTopology}).
\end{corollary}
\begin{proof}
Letting $f, g \in C_a(K)$ be any two functions, it is clear that $\min\{ f, g \}$ lies in every neighbourhood of both $f$ and $g$.
Similarly, $\max\{ f, g \} \subseteq \overline{\{ f \}} \cap \overline{ \{ g \} }$.
\end{proof}
%
\begin{corollary}
$\sigma_s( C(K), L )$ coincides with the topology of $C_p(K)$, i.e. that of pointwise convergence.
\end{corollary}
%
It is well-known that $C_p(K)$ is K-analytic and thus Lindelöf.
%
\begin{corollary}
$C_a(K)$ is Lindelöf.
\end{corollary}
%
\begin{corollary}
\label{cor:KToKHatHomeomorphism}
The mapping $K \to \widehat{K}: x \mapsto \delta_x$ is a homeomorphism when $\widehat{K}$ is equipped with $\sigma(L, C(K))$-topology.
\end{corollary}
\begin{proof}
By \zcref{prop:WeakStarTop}, the $\sigma(L, C(K))$-topology on $\widehat{K}$ coincides with the $\sigma(\widetilde{L}, C(K))$-topology.
Hence, the claim follows from \cite[Lemma 3.55]{src:Fabian:BanachSpaceTheory}.
\end{proof}
%
\begin{corollary}
\label{cor:SpecialisationOrderInCaKCoincidesWithFunctionOrder}
For $f, g \in C_a(K)$, $f \le g$ in the sense of \zcref{prop:WeakTop} if and only if $f \le g$ as functions.
\end{corollary}
%
An important property of the space $C_p(K)$ is its angelicity.
In particular, the class of (relatively) countably compact subsets coincides with the classes of (relatively) compact respectively (relatively) sequentially compact subsets.
A similar property holds for $C_a(K)$.
%
\begin{lemma}
\label{lem:CaKRelativelyCountablyCompactSubnetBounded}
Let $A \subseteq C_a(K)$ be relatively countably compact.
Then, for every net $(f_\alpha)_{\alpha \in I}$ in $A$ there is a cofinal, directed subset $J \subseteq I$  such that
%
\begin{equation}
\sup_{x \in K} \limsup_{\beta \in J} f_\beta \left( x \right) < \infty \, .
\end{equation}
\end{lemma}
\begin{proof}
Suppose the claim is false and fix some $x_1 \in K$ with $\limsup_\alpha f_\alpha(x_1) > 1$.
Then $J_1 = \{ \alpha \in I : f_\alpha(x_1) \ge 1 \}$ is cofinal and directed and $(f_\beta)_{\beta \in J_1}$ a subnet with $f_\beta(x_1) \ge 1$ for every $\beta \in J_1$.
By induction, there is a sequence $(x_n)_{n \in \mathbb{N}}$ in $K$ and for every $n \in \mathbb{N}$ a cofinal and directed set $J_n \subseteq I$ such that $g_\beta(x_m) \ge m$ for every $\beta \in J_n$ and $m \in \{ 1, \dots, n \}$.
Now, for every $n \in \mathbb{N}$, pick some $\alpha(n) \in J_n$ and consider the sequence $(f_{\alpha(n)})_{n \in \mathbb{N}}$.
By assumption, it has a cluster point $f \in C_a(K)$ and thus $f( x_m ) \ge m$ for every $m \in \mathbb{N}$.
Then, $\sup_{x \in K} f(x) = \infty$ contradicting the fact that the continuous function $f$ is bounded on the compact set $K$.
\end{proof}
%
\begin{corollary}
\label{cor:RelativeCompactnessEquivalences}
Let $A \subseteq C_a(K)$.
Then the following are equivalent:
%
\begin{enumerate}
\item $A$ is relatively compact
\item $A$ is relatively countably compact
\item $A$ is relatively sequentially compact
\end{enumerate}
\end{corollary}
\begin{proof}
$(i) \implies (ii)$ and $(iii) \implies (ii)$ are obvious.

$(ii) \implies (i)$:
Let $(f_\alpha)_{\alpha \in I}$ be a net in $A$.
By \zcref{lem:CaKRelativelyCountablyCompactSubnetBounded}, there is some $C \in \mathbb{R}$ and a subnet $(g_\beta)_{\beta \in J}$ with
%
\begin{equation}
\sup_{x \in K} \limsup_\beta g_\beta \left( x \right) \le C \, .
\end{equation}
%
Hence, $(g_\beta)_{\beta \in J}$ converges to the constant function with value $C$ on $K$.

$(ii) \implies (iii)$:
Let $(f_n)_{n \in \mathbb{N}}$ be a sequence in $A$.
By \zcref{lem:CaKRelativelyCountablyCompactSubnetBounded}, there is some $C \in \mathbb{R}$ and a subsequence $(g_n)_{n \in \mathbb{N}}$ with
%
\begin{equation}
\sup_{x \in K} \limsup_{n \to \infty} g_n \left( x \right) \le C \, .
\end{equation}
%
Hence, $(g_n)_{n \in \mathbb{N}}$ converges to the constant function with value $C$ on $K$.
\end{proof}
%
Consequently, $C_a(K)$ satisfies \zcref{def:angelicity:RelativeCompactnessEquivalence} of \zcref{def:angelicity}.
One may ask whether $C_a(K)$ is angelic as well, that is whether for every relatively countably compact subset $A \subseteq C_a(K)$ and every $\bar{a} \in \mathrm{cl}\, A$, there is a sequence in $A$ that converges to $\bar{a}$.
Unfortunately, in general, the answer turns out to be negative.
%
\begin{theorem}
\label{thm:Ca01IsNotAngelic}
$C_a([0,1])$ is not angelic.
\end{theorem}
\begin{proof}
Following \cite[Lemma II.3.5]{src:Arkhangelskii1:TopologicalFunctionSpaces}, there is a countable set $Z \subseteq C_p([0,1])$ such that for each $f \in Z$, $f([0,1]) \subseteq [0,1]$ and $0 \in \mathrm{cl}\, (Z \setminus \{ 0 \})$ but for which no sequence in $Z \setminus \{ 0 \}$ converges to $0$.
Since the topology of $C_a([0,1])$ is coarser than that of $C_p([0,1])$, $0$ is also in the $C_a(K)$-closure of $Z$.
Moreover, since all functions in $Z$ are nonnegative, a net in $Z$ $C_a([0,1])$-converges to $0$ precisely if it $C_p([0,1])$-converges to $0$.
At the same time, $Z$ is trivially seen to be relatively compact in $C_a(K)$, because every net in $Z$ $C_a(K)$-converges to the constant function with value one.
\end{proof}
%
However, the countable tightness survives the passage from $C_p(K)$ to $C_a(K)$.
%

%
\begin{corollary}
The space $C_a(K)$ has countable tightness.
\end{corollary}
\begin{proof}
For each $m, n \in \mathbb{N}$, the set
%
\begin{equation}
L_{m,n} = \left\{ \sum_{i = 1}^m t_i \delta_{x_i} \middle| t_1, \dots, t_m \in [0, \infty), x_1, \dots x_m \in K : \sum_{i = 1}^m t_i \le n \right\}
\end{equation}
%
is $\sigma(L,C_a(K))$-compact.
Since $L = \cup_{m,n \in \mathbb{N}} L_{m,n}$, the claim follows from \zcref{prop:DualIsSigmaCompactImpliesCountableTightness}.
\end{proof}
%
The failure of $C_a([0,1])$ to be angelic lies in the lack of the Fréchet-Urysohn property as was demonstrated in \zcref{thm:Ca01IsNotAngelic}.
In angelic spaces, that property implies the equivalence of the classes of compact, countably compact and sequentially compact subsets, a feature also lost in the asymmetric setting.
%
\begin{theorem}
Letting $\omega_1$ denote the first uncountable ordinal, there is a countably compact $A \subseteq C_a([0, \omega_1])$ that is not compact.
\end{theorem}
\begin{proof}
Set $A = \{ f_\alpha : \alpha < \omega_1 \}$, where for each ordinal $\alpha < \omega_1$,
%
\begin{equation}
f_\alpha : [0, \omega_1] \to \mathbb{R},
\beta \mapsto \begin{cases}
1 & \text{if } \beta \le \alpha \\
0 & \text{otherwise.}
\end{cases}
\end{equation}
%
Since the assignment $\alpha \mapsto f_\alpha$ is order-preserving and every countable set of ordinals has an upper bound less than $\omega_1$, it is obvious that every sequence in $A$ has a cluster point in $A$.
However, the net $(f_\alpha)_{\alpha < \omega_1}$ has no cluster point in $A$.
\end{proof}
%
The underlying reason why this equivalence fails, is that a countably compact subset of $C_a(K)$ is not necessarily $C_p(K)$-closed.
%
\begin{proposition}
\label{prop:CountablyCompactIsCompactCharacterisation}
Let $A \subseteq C_a(K)$ be countably compact and set $B = \mathrm{cl}_{C_p(K)} A$ with the subspace topology inherited from $C_a(K)$.
Then,
%
\begin{enumerate}
\item $A$ is precompact,
\item $A$ is relatively compact in $B$,
\item $A$ is compact if and only if $B \subseteq \downarrow A$,
\item $B$ is compact.
\end{enumerate}
\end{proposition}
\begin{proof}
This follows immediately from \zcref{lem:CountableCompactnessImpliesRelativeCompactnessInSymmetricClosure,thm:LindelöfCompactnessCondition,thm:LindelöfAndSymmetricClosureImpliesCompactness}.
\end{proof}
%
\zcref[S]{cor:RegularCountablyCompactImpliesSequentiallyCompact} shows that regularity and countable compactness implies sequential compactness.
However, as the following examples show, there are compact subsets of $C_a([0,1])$ that fail to be sequential.
%
\begin{example}
The set $Z$ from \zcref{thm:Ca01IsNotAngelic} together with the constant function with value one is compact and $T_0$ but not $T_1$.
It is however, not sequential, because no sequence in $(Z \cup \{ 1 \}) \setminus \{ 0 \}$ converges to $0$.
\end{example}
\begin{example}
\label{ex:CompactT1NonSequential}
To construct a compact $T_1$ set that is not sequential, let $(f_n)_{n \in \mathbb{N}}$ be an enumeration of $Z \setminus \{ 0 \}$.
Pick a continuous function $h : [0, 3] \to \mathbb{R}$ with $h([0,1]) = \{ 1 \}$, $h([1, 2]) = [0, 1]$ and $h((2, 3]) \subseteq (-\infty, 0)$.
Now, for every $n \in \mathbb{N}$ pick a continuous extension $g_n : [0, 3] \to \mathbb{R}$ of $f_n$ such that
%
\begin{itemize}
\item $g_n([1, 1 + 4^{-n}]) \subseteq [0, 1]$,
\item $g_n((1 + 4^{-n}, 1 + 2 \cdot 4^{-n})) \subseteq (1, \infty)$,
\item $g_n([1 + 2 \cdot 4^{-n}, 1 + 3 \cdot 4^{-n}]) \subseteq [0, 1]$,
\item $g_n((1 + 3 \cdot 4^{-n}, 1 + 4 \cdot 4^{-n})) \subseteq (-\infty, 0)$,
\item $g_n([1 + 4 \cdot 4^{-n}, 2]) = \{ 0 \}$
\item $g_n = h$ on $[2, 3]$
\end{itemize}
%
Then the set $D = \{ g_n : n \in \mathbb{N} \} \cup \{ h, 0 \} \subseteq C_a([0, 3])$ is countably compact, since $g_n$ converges to $h$ and thus compact, because it is countable.
Note, that $0 \in \mathrm{cl}\, \{ g_n : n \in \mathbb{N} \}$ because $g_n$ converges pointwise to zero on $(1, 2]$ and $g_n \le 0$ on $[2, 3]$.
By construction, $D$ is $T_1$ but not sequential, since no sequence in $D \setminus \{ 0 \}$ converges to $0$.
\end{example}
%
While perhaps not obvious on first sight, \zcref{ex:CompactT1NonSequential} is not Hausdorff.
This can be seen directly, by considering points in $(1, 2]$ that are close to $1$.
In fact, the existence of a compact Hausdorff space of countable tightness that is not sequential is independent of ZFC.
This amounts to the Moore-Mrówka problem (see e.g. \cite{src:Balogh:CompactHausdorffCountableTightness} and \cite[Classic Problem VI]{src:TopologyProceedings2:Problems}).
However, compact Hausdorff subspaces are regular and thus Fréchet-Urysohn by \zcref{cor:RegularCountablyCompactImpliesSequentiallyCompact}.

For an arbitrary set $Q$, let $\underline{C}(Q)$ denote the set of real-valued lower semicontinuous functions on $Q$ equipped with the subspace topology inherited from $( \mathbb{R}, u )^Q$.
Here, $u : \mathbb{R} \to \mathbb{R}$ is the asymmetric seminorm with $x \mapsto \max\{ x, 0 \}$ such that a net $(x_\alpha)_{\alpha \in I}$ in $(\mathbb{R}, u)$ converges to $x \in (\mathbb{R}, u)$ if and only if $\limsup_\alpha x_\alpha \le x$.
Consequently, a net $(f_\alpha)_{\alpha \in I}$ in $\underline{C}(Q)$ converges to some $f \in \underline{C}(Q)$ if and only if $\limsup_\alpha f_\alpha(q) \le f(q)$ for all $q \in Q$.
%
\begin{lemma}
\label{lem:AsymmetricHausdorffQuotient}
Let $B \subseteq C_a(K)$ be Hausdorff, $T \subseteq B$ and $\pi_2 : \widetilde{L} \to \widetilde{L} / \mathrm{ann}\, T$ the canonical quotient map where $\widetilde{L}$ is considered with the $\sigma(\widetilde{L}, C_a(K))$-topology.
Set $Q = \pi_2(\widehat{K})$ and define the map
%
\begin{equation}
\pi_1 : \mathrm{cl}_B T \to \underline{C}(Q), \text{ with } \left( \pi_1 f \right) \left( q \right)
=
\inf \left\{ f \left( k \right) \middle| k \in K : \delta_k \in \pi_2^{-1} \left( \left\{ q \right\} \right) \right\} \ ,
\end{equation}
%
for all $f \in \mathrm{cl}_B T$ and $q \in Q$.
Then,
%
\begin{enumerate}
\item \label{itm:ProjectionOnCoreEquality} For all $f \in T$ and all $k \in K$, $( \pi_1 f ) ( \pi_2 \delta_k ) = f( k )$,
\item for all $f \in T$, $\pi_1 f$ is a continuous function from $Q$ to $\mathbb{R}$,
\item $\pi_1$ is continuous and injective with a Hausdorff image,
\item \label{itm:CoreContinuity} A net $(f_\alpha)_{\alpha \in I}$ in $T$ converges to $f \in \mathrm{cl}_B T$ if and only if $(\pi_1 f_\alpha)_{\alpha \in I}$ converges to $\pi_1 f$,
\item If $B$ is regular, $\pi_1$ is a homeomorphism onto its image.
\end{enumerate}
\end{lemma}
\begin{proof}
First, note that $\mathrm{ann}\, T$ is $\sigma(\widetilde{L}, C_a(K))$-closed.
Hence, $\widetilde{L} / \mathrm{ann}\, T$ is Hausdorff and thus $Q$ is a compact, Hausdorff space.
By \zcref{cor:KToKHatHomeomorphism,prop:WeakStarTop}, the map
%
\begin{equation}
\psi : K \to Q, k \mapsto \pi_2 \left( \delta_k \right)
\end{equation}
%
is continuous and thus also closed and proper.
Using $\psi$ instead of $\pi_2$, one obtains
%
\begin{equation}
\label{eq:pi1fDefinitionViaPsi}
\left( \pi_1 f \right) \left( q \right)
=
\inf \left\{ f \left( k \right) \middle| k \in \psi^{-1} \left( \left\{ q \right\} \right) \right\} \ ,
\end{equation}
%
for all $f \in \mathrm{cl}_B T$ and all $q \in Q$.
Now, let $r \in \mathbb{R}$ and set $C = (\pi_1 f)^{-1}( (- \infty, r] )$ as well as $D = f^{-1}( (-\infty, r] )$.
Since $\psi$ is proper and $f$ is continuous, the infimum in \zcref{eq:pi1fDefinitionViaPsi} is always attained.
Hence, $q \in C$ implies that there is some $k \in K$ with $\psi(k) = q$ and $f(k) \le r$, i.e. $C \subseteq \psi(D)$.
The converse inequality is obvious, such that $C = \psi(D)$.
Since $D$ is closed, so is $C$ such that $\pi_1 f$ is lower semicontinuous and $\pi_1$ is well-defined.

$(i)$:
If $f \in T$ and $k, k' \in K$ with $\psi(k) = \psi(k')$, we have
%
\begin{equation}
\label{eq:CoreEqualityOnEquivalentElements}
f \left( k \right)
=
\left< f, \delta_k \right>
=
\left< f, \delta_{k'} \right>
=
f \left( k' \right) \, ,
\end{equation}
%
by definition such that \zcref{itm:ProjectionOnCoreEquality} follows.

$(ii)$:
Fix some $f \in T$, $r \in \mathbb{R}$ and let $U = (\pi_1 f)^{-1}( (-\infty, r) )$ as well as $V = f^{-1}( (-\infty, r ) )$.
Obviously, $\psi(V) = U$ such that $U$ is open if $\psi^{-1}( \psi(V) )$ is open.
However, by \zcref{eq:CoreEqualityOnEquivalentElements}, $V = \psi^{-1}( \psi(V) )$.
Since $f$ is continuous, $V$ is open and thus $U$ is open.
Consequently, $\pi_1 f$ is upper semicontinuous and thus continuous.

$(iii)$.
Let $(f_\alpha)_{\alpha \in I}$ be a net in $B$ that converges to some $f \in B$.
Then,
%
\begin{equation}
\begin{aligned}
\limsup_{\alpha} \left( \pi_1 f_\alpha \right) \left( q \right)
&=
\limsup_{\alpha} \inf_{k \in \psi^{-1} ( \{ q \} )} f_\alpha \left( k \right) \\
&\le
\inf_{k \in \psi^{-1} ( \{ q \} )} \limsup_{\alpha} f_\alpha \left( k \right)
\le
\inf_{k \in \psi^{-1} ( \{ q \} )} f \left( k \right)
=
\left( \pi_1 f \right) \left( q \right) \, ,
\end{aligned}
\end{equation}
%
for all $q \in Q$.
Now, let $f, g \in \mathrm{cl}_B T$ and set $a = \pi_1 f$ as well as $b = \pi_1 g$.
Since $\pi_1$ is continuous, $U \cap \pi_1(T) \neq \emptyset$ for all open $U \subseteq \underline{C}(Q)$ with $U \cap \pi_1(\mathrm{cl}_B T) \neq \emptyset$.
In particular, for all open neighbourhoods $U$ of $a$ and $V$ of $b$ with $U \cap V \neq \emptyset$, $U \cap V$ is open and $U \cap V \cap \pi_1(T) \neq \emptyset$.
Hence, if $a$ and $b$ have no disjoint neighbourhoods, there is a net $(f_\alpha)_{\alpha \in I}$ in $T$ such that $(\pi_1 f_\alpha)_{\alpha \in I}$ converges to both $a$ and $b$.
However, then
%
\begin{equation}
f \left( k \right)
\ge
a \left( \psi(k) \right)
\ge
\limsup_\alpha \left( \pi_1 f_\alpha \right) \left( \psi(k) \right)
=
\limsup_\alpha f_\alpha \left( k \right) \, ,
\end{equation}
%
for all $k \in K$.
Hence, $(f_\alpha)_{\alpha \in I}$ converges to $f$ and similarly, to $g$ as well.
Since $\mathrm{cl}_B T$ is Hausdorff, it follows that $f = g$.
Consequently, $\pi_1$ is injective and has a Hausdorff image.

$(iv)$:
The forward implication is simply the continuity of $\pi_1$.
For the converse, let $f \in \mathrm{cl}_B T$ and $(f_\alpha)_{\alpha \in I}$ be a net in $T$ such that $(\pi_1 f_\alpha)_{\alpha \in I}$ converges to $\pi_1 f$.
Then, by applying \zcref{itm:ProjectionOnCoreEquality} for all $k \in K$,
%
\begin{equation}
\limsup_{\alpha} f_\alpha \left( k \right)
=
\limsup_{\alpha} \left( \pi_1 f_\alpha \right) \left( \psi \left( k \right) \right)
\le
\left( \pi_1 f \right) \left( \psi \left( k \right) \right)
\le
f \left( k \right) \, .
\end{equation}

$(v)$:
Consider the function $g = \pi_1^{-1} \upharpoonright \pi_1(T)$.
Applying \cite[p. 81, Theorem 1]{src:Bourbaki:GeneralTopologyI}, $g$ has a unique continuous extension to $\pi_1(\mathrm{cl}_B T)$ which by definition coincides with $\pi_1^{-1}$.
\end{proof}
%
\begin{theorem}
Let $B \subseteq C_a(K)$ be countably compact and regular.
Then $B$ is Fréchet-Urysohn.
\end{theorem}
\begin{proof}
Let $A \subseteq B$ and $\bar{a} \in \mathrm{cl}_B A$.
Since $C_a(K)$ is countably tight, there is a sequence $(a_n)_{n \in \mathbb{N}}$ in $A$ with $\bar{a}$ as a cluster point.
Setting $T = \{ a_n : n \in \mathbb{N} \} \cup \{ \bar{a} \}$ and applying \zcref{lem:AsymmetricHausdorffQuotient}, we have the canonical projection $\pi_2 : \widetilde{L} \to \widetilde{L} / \mathrm{ann}\, T$, $Q = \pi_2(\widehat{K})$ and the continuous, injective function $\pi_1 : \mathrm{cl}_B T \to \underline{C}(Q)$ with a Hausdorff image.

Clearly, $(\mathrm{span}\, T, \mathrm{cone}\, Q)$ is a duality and the topology on $Q$ is given by the subspace topology of $\sigma( \mathrm{cone}\, Q, \mathrm{span}\, T)$.
Since $T$ is countable, \zcref{lem:CountableBasisImpliesMetrisableCompactaInDual} applies such that $Q$ is metrisable and thus, separable.
Consequently, there is a dense sequence $(q_n)_{n \in \mathbb{N}}$ in $Q$.
Now, writing $\bar{g} = \pi_1(\bar{a})$ as well as $g_n = \pi_1(a_n)$, we proceed by showing that the pseudocharacter of $C$ at $\bar{g}$ is countable.
Indeed, for every $n \in \mathbb{N}$, consider the open neighbourhood of $\bar{g}$ given by
%
\begin{equation}
U_n = \left\{ g \in C \middle| \forall m \in \{ 1, \dots, n \} : g \left( q_m \right) < \bar{g} \left( q_m \right) + \frac{1}{n} \right\} \, .
\end{equation}
%
If $g \in \cap_{n \in \mathbb{N}} U_n$, it follows that $g(q_m) \le \bar{g}(q_m)$ for all $m \in \mathbb{N}$.
Moreover, recalling from \zcref{lem:AsymmetricHausdorffQuotient} that $\bar{g}$ is continuous and $g$ is lower semicontinuous, it follows that $g \le \bar{g}$.
However, since $C$ is Hausdorff, \zcref{cor:SpecialisationOrderInCaKCoincidesWithFunctionOrder} implies $g = \bar{g}$.
Because $B$ is regular, $\pi_1$ is a homeomorphism onto $C$ such that $C$ is regular as well and the claim follows from \zcref{lem:CountabilityAndPseudocharacter}.
\end{proof}
%
\begin{corollary}
\label{cor:RegularCountablyCompactImpliesSequentiallyCompact}
Every regular, countably compact subset of $C_a(K)$ is sequentially compact.
\end{corollary}
%
\begin{corollary}
For a Hausdorff subset $A$ of $C_a(K)$, the following are equivalent:
%
\begin{enumerate}
\item $A$ is compact
\item $A$ is countably compact and $\mathrm{cl}_{C_p(K)} A \subseteq \downarrow A$,
\item $A$ is sequentially compact and $\mathrm{cl}_{C_p(K)} A \subseteq \downarrow A$.
\end{enumerate}
%
Moreover, in the affirmative case, $A$ is Fréchet-Urysohn and thus angelic.
\end{corollary}
\begin{proof}
Apply \zcref{prop:CountablyCompactIsCompactCharacterisation} and the fact that compact Hausdorff spaces are regular.
\end{proof}
%
\section{The Weak Topology on an Asymmetrically Normed Space}
\label{sec:WeakTopologyOnAsymmetricallyNormedSpace}
%
It is well-known that the weak topology on a normed space is angelic \cite[Proposition 3.108]{src:Fabian:BanachSpaceTheory}.
In this section we apply the results from the preceeding \zcref[noref]{sec:CaK} to the setting of asymmetrically normed spaces.
%
\begin{proposition}
Let $(X, p)$ be an asymmetrically normed space.
Then, $(X, (X, p)^*)$ is homeomorphic to a subset of $C_a(K)$ where
%
\begin{equation}
K = \left\{ \phi \in Y : \sup_{p(x) \le 1} \phi \left( x \right) \le 1 \right\}
\end{equation}
%
is compact in the $\sigma(Y, X)$ topology.
\end{proposition}
\begin{proof}
That $K$ is compact is well-known \cite[Theorem 2.4.3]{src:Cobzaş:FunctionalAnalysisinAsymmetricNormedSpaces}.
Let $f : (X, (X, p)^*) \to C_a(K)$ be given by $f(x)(\phi) = \phi(x)$ for all $x \in X$ and $\phi \in K$.
Then, $f$ is a homeomorphism by \zcref{prop:WeakStarTop,cor:XandXDualAreInDuality,cor:CaKConvergence}.
\end{proof}
%
\begin{corollary}
An asymmetrically normed space has countable tightness in its weak topology.
\end{corollary}
%
\begin{corollary}
Let $X$ be an asymmetrically normed space and $A \subseteq X$ a subset equipped with the weak topology.
If $A$ is countably compact and $\sigma_s(X, X^*)$-closed in $X$, then $A$ is compact.
If, in addition, $A$ is Hausdorff, then it is Fréchet-Urysohn and angelic.
\end{corollary}
%
\appendix
\printbibliography
\end{document}
